% Book pp. 122-125 (PDF pp. 123-126)
\section{Complex Conjugates Theorem}

\textbf{Complex Numbers}: A complex number is expressed in the form $a + bi$, where a
and b are real numbers, and $i$ is the imaginary unit defined as $i^2 = -1$, or
$i = \sqrt{-1}$

\textbf{Roots of Polynomials}: When dealing with polynomial equations, the Fundamental
Theorem of Algebra states that every non-constant polynomial has at least one
complex root. This means that if a polynomial has real coefficients, any non-real
roots must occur in conjugate pairs.

If $a + bi$ (where b $\ne$ 0) is a root of a polynomial, then its conjugate $a - bi$ is
also a root.

\textbf{Example}: For a polynomial like $x^2 + 4$:

The roots can be found by rearranging to $x^2 = -4$.

Taking the square root gives $x = \pm 2i$

Here, $2i$ and $-2i$ are complex roots that come in conjugate pairs.

\textbf{Note that:}

\textbf{Complex roots} always come in \textbf{conjugate pairs}.

If $a + bi$ is a root of a polynomial, then $a - bi$ is also a root.

This property helps us find all the roots of polynomials, ensuring we account for
both real and complex solutions.

The \textbf{Complex Conjugate Theorem} states that if a polynomial has a complex
root, its conjugate must also be a root. This is an important concept in algebra that
helps us understand the nature of polynomial equations and their solutions.

Let us go through the following examples to explore how to use the \textbf{Complex
conjugate theorem.}

For example, if we are given one complex factor of a polynomial, we know
another factor must be the complex conjugate. This gives us two factors. We can
then form a quadratic equation and use algebraic division to find another factor.
Or we may be given two real factors, for example, $x = 2$ and $x = -3$. This means
that $x - 2 = 0$ or $x + 3 = 0$

$\therefore (x - 2)(x + 3) = 0$

$x^2 + 3x - 2x - 6 = 0$

$x^2 + x - 6 = 0$

\begin{example}{Example 3.15}
Find all the roots of $f(x) = x^3 + 8x^2 + 21x + 20$ if one root is -2 $-$ $i$.

\solution
$f(x) = x^3 + 8x^2 + 21x + 20$

Since the factor given $(-2 - i)$ is complex in nature, we need the conjugate as
well: $-2 + i$

$[x - (-2 - i)][(x - (-2 + i)]$

$[x + 2 + i][x + 2 - i]$

$(x + 2)^2 - (i)^2$

$x^2 + 4x + 4 - i^2$

NB: $-i^2 = 1$

$x^2 + 4x + 4 - (-1)$

$x^2 + 4x + 5$

We can use long division to find the other factors:
\[
\begin{array}{r@{\;}r}
 & x + 4 \\ \cline{2-2}
x^2 + 4x + 5\,\big) & x^3 + 8x^2 + 21x + 20 \\
 & \underline{-(x^3 + 4x^2 + 5x + 20)} \\
 & 4x^2 + 16x + 20 \\
 & \underline{-(4x^2 + 16x + 20)} \\
 & 0
\end{array}
\]
$x + 4 = 0$

$x = -4$

Therefore, the roots of $(x)$ are $-4$, $-2 + i$ and $-2$ -- $i$
\end{example}

\begin{example}{Example 3.16}
Find all the roots of $f(x) = 2x^3 - 13x^2 + 26x - 10$ if one root is $3 + i$.

\solution
$f(x) = 2x^3 - 13x^2 + 26x - 10$

Since the factor given $3 + i$. is complex in nature, we need the conjugate as well:
$3 - i$

$[x - (3 + i)][(x - (3 - i)]$

$[x - 3 - i][x - 3 + i]$

$(x - 3)^2 - (i)^2$

$x^2 - 6x + 9 - i^2$

NB: $-i^2 = -1$
\booknote[S3-14]{$-i^2 = 1$, as Example 3.15 states; the next line correctly
subtracts $(-1)$.}

$x^2 - 6x + 9 - (-1)$

$x^2 - 6x + 10$

We can use long division to find the other factor.
\[
\begin{array}{r@{\;}r}
 & 2x - 1 \\ \cline{2-2}
x^2 - 6x + 10\,\big) & 2x^3 - 13x^2 + 26x - 10 \\
 & \underline{-(2x^3 - 12x^2 + 20x + 0)} \\
 & -x^2 + 6x - 10 \\
 & \underline{-(-x^2 + 6x - 10)} \\
 & 0
\end{array}
\]
$2x - 1 = 0$

$x = \frac{1}{2}$

Therefore, the roots of $(x)$ are $= \frac{1}{2}$, $3 + i$ and $3$ -- $i$
\end{example}
